\section{A Core Calculus for Transparent Outputs}
\label{sec:core}

We start by introducing a core language for supporting transparent outputs, which is the basis of \OurLang.
The term syntax (\secref{core:syntax}) is uncontroversial: it is pure and untyped, and provides datatypes,
records and matrices (although we omit a presentation of matrices here). The syntax of values is less
standard: each part of a value has a unique \emph{address} $\alpha$ which serves to identify that part of the
value in a dependence graph. We give a big-step operational semantics that evaluates a term to both a value
and a \emph{dynamic dependence graph} for that value (\secref{core:semantics}), capturing how parts of that
value depend on parts of the input. The core language is fairly expressive, supporting deep pattern-matching
and mutual recursion; however our implementation also provides a surface language that desugars into the core,
providing piecewise function definitions, list comprehensions and other conveniences, as shown in
\figrefTwo{example:moving-average:source}{example:non-renewables:source}, which would otherwise complicate the
core.

\paragraph{Some notation.} We write to ${\seq{x}}$ denote a finite sequence of elements $\seqRange{x_1}{x_n}$,
with $\seqEmpty$ as the empty sequence. Concatenation of sequences is written $\seq{x} \concat \seq{x}'$; we
also write $x \cons \seq{x}$ for cons (prepend) and $\seq{x} \cons x$ for snoc (append). We write
$\set{\seq{\bind{k}{x}}}$ to denote a finite map, i.e.~a set of pairs
$\seqRange{\bind{k_1}{x_1}}{\bind{k_n}{x_n}}$ where keys $k_i$ are pairwise unique. If $X$ and $Y$ are sets,
we write $X \disjunion Y$ to mean $X \cup Y$ where $X$ and $Y$ are disjoint, and also $x \cdot X$ or $X \cdot
x$ to mean $X \disjunion \set{x}$. If $X$ and $Y$ happen to be finite maps, we further constrain $X \disjunion
Y$ to the case where $\dom{X}$ and $\dom{Y}$ are disjoint.

\subsection{Syntax}
\label{sec:core:syntax}

\begin{figure}
   \begin{minipage}[t]{0.48\textwidth}
      \begin{tabularx}{\textwidth}{rL{2.7cm}L{3cm}}
         &\textbfit{Core term}&
         \lowlight{$\Expr\;\Gamma\;A$}
         \\
         $e ::=$
         &
         $\exSugarNew{\raw{s}}{e}$
         &
         desugared term
         \\
         &
         $\exInt{n}$
         &
         integer
         \\
         &
         $\exStr{s}$
         &
         string
         \\
         &
         $\exVar{x}$
         &
         variable
         \\
         &
         $f$
         &
         external function
         \\
         &
         $\exRec{\seq{\bind{x}{e}}}$
         &
         record
         \\
         &
         $\exRecProj{e}{x}$
         &
         record projection
         \\
         &
         $\exDict{\seq{\bind{e}{e}}'}$
         &
         dictionary
         \\
         &
         $\exConstr{c}{e}$
         &
         data value
         \\
         &
         $\exLambda{\upsigma}$
         &
         anonymous function
         \\
         &
         $\exApp{e}{e'}$
         &
         application
         \\
         &
         $\exLet{x}{e}{e'}$
         &
         let
         \\
         &
         $\exLetRec{\seq{\bind{x}{\upsigma}}}{e}$
         &
         recursive let
      \end{tabularx}
   \end{minipage}\\[3mm]
\end{figure}


\subsubsection{Terms}
\label{sec:core:expressions}

\figref{syntax} defines the expressions $e$ of the language, which include variables $x$, integer constants
$n$, let-bindings $\exLet{x}{e}{e'}$, record construction $\exRec{\seq{\bind{x}{e}}}$ and record projection
$\exRecProj{e}{x}$. We also have (saturated) constructor expressions $\exConstr{c}{\seq{e}}$, where $c$ ranges
over data constructors; the language is parameterised by a finite map $\Sigma$ from constructors $c$ to
constructor arities $\datatype{c} \in \mathbb{N}$. Function application comes in two forms: the usual
$\exApp{e}{e'}$, and (saturated) \emph{foreign function applications} $\exForeignApp{f}{\seq{e}}$ for calling
foreign functions (\secref{core:foreign} below). The final two expression forms are anonymous functions
$\exFun{\sigma}$ and mutually recursive let-bindings $\exLetRecElim{\seq{x : \sigma}}{e}$, where $\sigma$ is a
pattern-matching construct called an \emph{eliminator} (\secref{core:continuations} below).

\subsubsection{Foreign functions}
\label{sec:core:foreign}
The language is also parameterised by a finite map $\Phi$ from variables $f$ to arities $\Phi(f) \in \Nat$ of
the foreign function they denote. For the graph semantics in \secref{core:semantics} below, every foreign
function $f$ is required to provide an interpretation $\interpret{f}$ that returns both the result of calling
$f$ and a dependence graph capturing how the result depends on the arguments supplied to $f$.

\subsubsection{Continuations and eliminators}
\label{sec:core:continuations}
A \emph{continuation} $\kappa$  is either a term $e$ or a eliminator $\sigma$, and describes how an execution
proceeds after a value is pattern-matched. An \emph{eliminator} is a deep pattern-matching construct based on
tries~\cite{hinze00,peytonjones2021}; for a language with rich structured values like records and data types,
we find this makes for a cleaner presentation than a single shallow elimination form for each type. The
piecewise function definitions of the surface language desugar into eliminators.

An eliminator specifies how values of a particular shape are matched, and for a given value, determines how
any pattern variables get bound and the continuation $\kappa$ which will be executed under those bindings. A
variable eliminator $\elimVar{x}{\kappa}$ says how to match any value (as variable $x$) and continue as
$\kappa$. A record eliminator $\elimRecord{\seq{x}}{\kappa}$ says how to match a record with fields $\seq{x}$,
and provides a continuation $\kappa$ for sequentially matching the values of those fields. Lastly, a
constructor eliminator $\elimConstr{\seq{c \mapsto \kappa}}$ provides a branch $c \mapsto \kappa$ for each
constructor $c$ in $\seq{c}$, where each $\kappa$ specifies how any arguments to $c$ will be matched. (We
assume that any constructors matched by an eliminator all belong to the same data type, but this is not
enforced in the core language.)

\subsubsection{Values and environments}
\label{sec:core:values}

We define addressed \emph{values} $v$ mutually inductively with \emph{raw values} $\raw{v}$, \todo{sentences
separated by comma} a raw value does not have an associated address. These addresses, explained in
\secref{core:addresses}, will (later) \todo{weird forward-reference} form the vertices of the graphs that we
build in the graph semantics.

Raw values $\raw{v}$ include integers $n, m$; records  $\exRecord{\seq{x \mapsto v}}$; and constructor values
$\exConstr{c}{\seq{v}}$. We also define function closures as $\exClosure{\gamma}{\rho}{\sigma}$ where $\sigma$
is the function body, $\gamma$ is its environment, and to support mutual recursion, ${\rho}$ is the (possibly
empty) set of named functions with which $\sigma$ was mutually defined. \emph{Environments} are finite maps
from variables to values. Because foreign functions in the core language are not first-class and calls
$\exForeignApp{f}{\seq{e}}$ are saturated, i.e.~$\length{\seq{e}} = \Phi{(f)}$, the surface language provides
a top-level environment which maps every foreign function name ${f}$ of arity $n$ to the closure
$\annClosure{\envEmpty}{\envEmpty}{\elimVar{x_1}{\elimVar{\iter}{\elimVar{x_{n}}{\exForeignApp{f}{\seq{x}}}}}}{\alpha}$,
with $\alpha$ fresh, emulating first-class foreign functions.

\subsubsection{Addresses for partial values}
\label{sec:core:addresses}
\todo{Introduce addresses early on.} The addresses $\alpha$ that decorate values act as identifiers, uniquely
corresponding those values to specific vertices in a dependence graph. There are many sensible ways of
assigning addresses (our particular implementation uses an integer counter to generate new ones) as long as
some freshness condition is satisfied -- ensuring that there are no unwanted address collisions in a given
value i.e.~graph \todo{``value i.e.~graph''? what would a ``wanted'' address collision be?}. For simplicity,
we elide this side-condition \todo{what side condition?} in the graph semantics (next) and assume it is met
whenever a new address $\alpha$ is introduced. \todo{This paragraph unclear.}

% The point of the disjointness conditions on graph union was intended to achieve this, and does at least
% ensure that graph vertices don’t collide unexpectedly. However, I’m not sure that’s sufficient to avoid
% problems. For example, there doesn’t seem to be anything in the closeDefs definition that prevents a
% newly-created closure being allocated an address that is already in use somewhere in γ. This probably needs
% a mention and isn’t likely to be easy to formalise (a side-condition which says “α_i fresh for all i” isn’t
% much better than just saying this in plain English outside of the definition).

\subsection{Operational Semantics}
\label{sec:core:semantics}

We now give a big-step operational semantics for the core language, which evaluates terms to values and builds
a dynamic dependence graph for any such value. First we introduce the idea of a dynamic dependence graph.

\subsubsection{Dynamic Dependence Graphs}
\label{sec:core:dependence-graphs}

A \emph{dynamic dependence graph}~\cite{agrawal90} (hereafter \emph{dependence graph}) is a directed acyclic
graph $G = (V,E)$ with a set $V$ of \emph{vertices} and a set $E \subseteq V \times V$ of edges. In the
dependence graph for a particular program, vertices $\alpha, \beta \in V$ represent \emph{parts} of values
(either supplied to the program as inputs or produced during evaluation) and edges $(\alpha, \beta) \in E$
indicate that, in the evaluation of the program, the partial value associated to $\beta$ \textit{depends on}
on the partial value associated to $\alpha$. This diverges somewhat from traditional approaches to dynamic
dependence graphs in only considering one type of edge (data dependency) rather than separate data and control
dependencies; moreover our edges point in the direction of dependent vertices, whereas in the literature the
other direction is somewhat more common.

During evaluation, the output graph $G$ is built by specifying new graph fragments and then combining them;
the following notation is convenient for specifying a graph fragment where a fresh vertex $\alpha$ depends on
a set $V$ of preexisting vertices.

\begin{definition}[In-star notation]
Write $\inStar{V}{\alpha}$ as shorthand for the star graph $(\set{\alpha} \disjunion V, V \times
\set{\alpha})$.
\end{definition}

\todo{Illustrate dependence graphs, perhaps using moving average example?}

\begin{figure}
    {\small \flushleft \shadebox{$\seq{v}, \upkappa \matchFwdS \upgamma, \upkappa', V$}\hfill \textit{with }$\dom{\upgamma} = \bigcup(\seq{\BV(w)})$%
    \begin{smathpar}
       \inferrule*[
          lab={\ruleName{$\matchFwdS$-seq-empty}}
       ]
       {
          \strut
       }
       {
          \seqEmpty, \upkappa \matchFwdS \varnothing, \upkappa, \varnothing
       }
       %
       \and
       %
       \inferrule*[
          lab={\ruleName{$\matchFwdS$-seq-non-empty}}
       ]
       {
          v, \upsigma \matchFwdS \upgamma, \upkappa, V
          \\
          \seq{u}, \upkappa \matchFwdS \upgamma', \upkappa', V'
       }
       {
          v \cons \seq{u},
          \upsigma
          \matchFwdS
          \upgamma \disjunion \upgamma',
          \upkappa',
          V \cup V'
       }
    \end{smathpar}}

    {\small \flushleft \shadebox{$v, \upsigma \matchFwdS \upgamma, \upkappa, V$}\hfill \textit{with }$\dom{\upgamma} = \BV(w)$%
    \begin{smathpar}
       \inferrule*[lab={\ruleName{$\matchFwdS$-var}}]
       {
          \strut
       }
       {
          v, \elimVar{x}{\upkappa} \matchFwdS \set{\bind{x}{v}}, \upkappa, \varnothing
       }
       %
       \and
       %
       \inferrule*[
          lab={\ruleName{$\matchFwdS$-record}}
       ]
       {
          \seq{u}, \upkappa \matchFwdS \upgamma, \upkappa', V
          \\
          \exRec{\seq{\bind{y}{u}}} \subseteq \exRec{\seq{\bind{x}{v}}}
       }
       {
          \annRec{\seq{\bind{x}{v}}}{\upalpha},
          \elimRec{\seq{y}}{\upkappa}
          \matchFwdS
          \upgamma, \upkappa', V \cup \set{\upalpha}
       }
       %
       \and
       %
       \inferrule*[lab={\ruleName{$\matchFwdS$-constr}}]
       {
          \seq{v}, \upkappa' \matchFwdS \upgamma, \upkappa^\twoPrime, V
       }
       {
          \annConstr{c'}{\seq{v}}{\upalpha}, \elimConstr{\elimBind{c'}{\upkappa'} \cons \seq{\elimBind{c}{\upkappa}}}
          \matchFwdS
          \upgamma, \upkappa^\twoPrime, V \cup \set{\upalpha}
       }
    \end{smathpar}}
 \caption{Core semantics: matching}
 \end{figure}

\begin{figure}
   {\small \flushleft \shadebox{$\gamma, \seq{e} \evalSugS \seq{v}$}%
   \begin{smathpar}
      \inferrule*[
         lab={\ruleName{$\evalSugS$-seq}}
      ]
      {
         \gamma, e_i \evalSugS v_i
         \quad
         (\forall i \numleq \length{\seq{e}})
      }
      {
         \gamma, \seq{e} \evalSugS \seq{v}
      }
   \end{smathpar}}
   {\small \flushleft \shadebox{$\gamma, e \evalSugS v$}%
   \begin{smathpar}
      \inferrule*[
      lab={\ruleName{$\evalSugS$-var}}
      ]
      {
         \strut
      }
      {
         \gamma \cons (\bind{x}{v}), \exVar{x} \evalSugS v
      }
      %
      \and
      %
      \inferrule*[
         lab={\ruleName{$\evalSugS$-int}}
      ]
      {
         \strut
      }
      {
         \gamma, \exInt{n}
         \evalSugS
         \exInt{n}
      }
      %
      \and
      %
      \inferrule*[lab={\ruleName{$\evalSugS$-record}}]
      {
         \gamma, \seq{e} \evalSugS \seq{v}
      }
      {
         \gamma, \exRec{\seq{\bind{x}{e}}}
         \evalSugS
         \exRec{\seq{\bind{x}{v}}}
      }
      %
      \and
      %
      \inferrule*[
         lab={\ruleName{$\evalSugS$-project}}
      ]
      {
         \gamma, e \evalSugS \exRec{\seq{\bind{x}{v}} \cons (\bind{y}{u})}
      }
      {
         \gamma, \exRecProj{e}{y} \evalSugS u
      }
      %
      \and
      %
      \inferrule*[lab={\ruleName{$\evalSugS$-constr}}]
      {
         \gamma, \seq{e} \evalSugS \seq{v}
      }
      {
         \gamma, \exConstr{c}{\seq{e}}
         \evalSugS
         \exConstr{c}{\seq{v}}
      }
      %
      \and
      %
      \inferrule*[lab={\ruleName{$\evalSugS$-lambda}}]
      {
         \strut
      }
      {
         \gamma, \exFun{\sigma}
         \evalSugS
         \exClosure{\gamma}{\envEmpty}{\sigma}
      }
      %
      \and
      %
      \inferrule*[
         lab={\ruleName{$\evalSugS$-app}},
         width=3.3in,
      ]
      {
         \gamma, e \evalS \exClosure{\gamma_1}{\rho}{\sigma}
         \\
         \gamma, e' \evalS v
         \\
         \gamma_1, \rho \closeDefs \gamma_2
         \\
         v, \sigma \match \gamma_3, e^\dag
         \\
         \gamma_1 \concat \gamma_2 \concat \gamma_3, e^\dag \evalSugS u
      }
      {
         \gamma, \exApp{e}{e'}
         \evalSugS
         u
      }
      %
      \and
      %
      \inferrule*[
         lab={\ruleName{$\evalS$-foreign}}
      ]
      {
         \gamma, \seq{e} \evalSugS \seq{v}
         \\
         \interpret{f}(\seq{v}) = u
      }
      {
         \gamma, \exForeignApp{f}{\seq{e}}
         \evalSugS
         u
      }
      %
      \and
      %
      \inferrule*[
         lab={\ruleName{$\evalSugS$-let}}
      ]
      {
         \gamma, e \evalSugS v
         \\
         \gamma \cons (\bind{x}{v}), e' \evalSugS v'
      }
      {
         \gamma, \exLet{x}{e}{e'}
         \evalSugS
         v'
      }
      %
      \and
      %
      \inferrule*[
         lab={\ruleName{$\evalSugS$-let-rec}}
      ]
      {
         \gamma, \set{\seq{\bind{x}{\sigma}}} \closeDefs \gamma'
         \\
         \gamma \concat \gamma', e \evalSugS v
      }
      {
         \gamma, \exLetRec{\seq{\bind{x}{\sigma}}}{e}
         \evalSugS
         v
      }
   \end{smathpar}}
   \\[2mm]
   {\small \flushleft \shadebox{$\gamma, \rho \closeDefs \gamma'$}\hfill\textit{ with }$\dom{\rho} = \dom{\gamma}$%
   \begin{salign}
      \gamma, \rho
      &\closeDefs
      \set{\bind{x}{\exClosure{\gamma}{\rho}{\rho(x)}} \mid x \in \dom{\rho}}
   \end{salign}
   }
   \caption{Operational semantics: evaluation}
   \label{fig:core:eval}
\end{figure}

\begin{figure}
  {\small \flushleft \shadebox{$\gamma, \rho, V \closeDefs \gamma', G$}
   \hfill\textit{ with }$\dom{\gamma'} = \dom{\rho}$%
  \begin{smathpar}
     \inferrule*[]
     {
        v_i = \annClosure{\gamma}{\rho}{\rho(x_i)}{\alpha_i}
        \\
        G_i = \set{\bind{\alpha_i}{V}}
        \quad
        (\forall i \numleq \length{\seq{x}})
     }
     {
        \gamma, \rho, V
        \closeDefs
        \set{\seq{\bind{x}{v}}},
        \textstyle{\bigdisjunion}\seq{G}
     }
  \end{smathpar}
  }
  \vspace{-0.5cm}
  \caption{``Closing over'' recursive definitions}
  \label{fig:core:close-defs}
\end{figure}


\subsubsection{Pattern matching}
\label{sec:core:semantics:pattern-matching}

\figref{core:pattern-matching} defines the pattern-matching judgement $\seq{v}, \kappa \match \gamma, e, V$.
Rather than matching a single value $v$, the judgement matches a ``stack'' of values $\seq{v}$ against a
continuation $\kappa$, returning the selected branch $e$, an environment $\gamma$ providing bindings for the
free variables in $e$, and the set $V$ of addresses found in the matched portions of $\seq{v}$.

The \ruleName{$\match$-done} rule expresses exactly when using a \emph{term} $e$ as the continuation may
succeed: when the stack of values to be matched is empty, and so $e$ is treated as the selected term to be
executed. A variable eliminator $\elimVar{x}{\kappa}$ matches any value $v$ and and selects $\kappa$ as the
continuation; the empty set of addresses is returned, reflecting the fact that no part of $v$ was consumed.
The other two rules require an \emph{eliminator} $\sigma$ as the continuation, which is used to match the head
of a non-empty stack $v \cons \seq{v}$; any relevant subvalues are unpacked and pushed onto the tail $\seq{v}$
to be recursively matched by the continuation by $\sigma$. Specifically, the record eliminator
$\elimRecord{\seq{y}}{\kappa}$ matches a record of the form $\annot{\exRec{\seq{\bind{x}{v}}}}{\alpha}$ as
long as the variables in $\seq{y}$ are also fields in $\seq{x}$, returning the address $\alpha$ associated
with the record; this condition is specified by the premise $\exRec{\seq{\bind{y}{u}}} \subseteq
\exRec{\seq{\bind{x}{v}}}$, which simultaneously projects out the corresponding values of $\seq{u}$ from
$\seq{v}$. Lastly, a constructor eliminator with a branch $(\elimBind{c}{\kappa})$ matches any constructor
value of the form $\annot{\exConstr{c}{\seq{v}}}{\alpha}$, returning the address $\alpha$ associated with the
constructor.

\subsubsection{Evaluation}

\figref{core:eval} defines a big-step evaluation relation $\gamma, e, V \evalS v, G$, stating that term $e$,
under an environment $\gamma$ and vertex set $V$, evaluates to a value $v$ plus dependence graph~$G$. The
input set $V$ corresponds to the inputs consumed by the current active function call from which $e$ is being
evaluated, and can be thought of as the ``ambient'' set of  dependencies \todo{arising from what?}; this is
initially empty for a top-level term, and may grow whenever a function is applied to an argument and
pattern-matching occurs (i.e. in the \ruleName{$\evalS$-app} rule, which we explain later) \todo{and when it
is discarded?}.

The \ruleName{$\evalS$-var} rule for evaluating variables is standard, and the returned dependence graph is
empty; no new addresses are allocated as a result of simply looking up a variable. The
\ruleName{$\evalS$-record-project} is similar: if $e$ evalutes to a record $\annRec{\seq{\bind{x}{v}} \cons
(\bind{y}{u})}{\alpha}$, then $e.y$ evalutes to the value $u$ of field $y$, discarding address $\alpha$ of the
record. Most of the other evaluation rules can generally be understood in terms of two recurring patterns.

First, for any introduction rule that creates a new (partial) value, there is no choice \todo{really?} but to
assign that partial value a fresh address, which is then linked to each of the ambient vertices in $V$ in the
resulting graph. These cases include \ruleName{$\evalS$-int}, \ruleName{$\evalS$-function},
\ruleName{$\evalS$-record}, and \ruleName{$\evalS$-constr}. For example, \ruleName{$\evalS$-int} evaluates an
integer $n$ to its value form $n_{\alpha}$ assigned a fresh address $\alpha$; the resulting dependence graph
is then $\inStar{V}{\alpha}$, denoting that evaluating to $n_\alpha$ depended on any inputs that were found in
$V$. Likewise, \ruleName{$\evalS$-function} evaluates an anonymous function $\exFun{\sigma}$ by constructing
the closure $\exClosure{\gamma}{\envEmpty}{\sigma}_\alpha$ with fresh address $\alpha$; here, the closure
captures the current environment $\gamma$ and uses $\envEmpty$ as the set of mutually recursive definitions.

The second recurring pattern is that, for any rule which involves evaluating subterms of a term, the
corresponding subgraphs will have disjoint domains (meaning any addresses allocated for those subgraphs do not
overlap), and their union is the dependence graph for the overall term. This general principle is reflected in
the auxiliary evaluation relation $\gamma, \seq{e}, V \evalS \seq{v}, G$ (bottom of \figref{core:eval}), which
evalutes $e_i$ in a sequence of terms $\seq{e}$ to a value $v_i \in \seq{v}$ and graph $G_i \in \seq{G}$, and
returns the dependence graph $\textstyle{\bigdisjunion} \seq{G}$. We make use of this judgement in the rules
\ruleName{$\evalS$-record} for $\exRec{\seq{\bind{x}{e}}}$, \ruleName{$\evalS$-constr} for $c(\seq{e})$, and
\ruleName{$\evalS$-foreign-app} for $f(\seq{e})$, which all involve evaluating a sequence $\seq{e}$ of
subterms; in the last rule, $\hat{f}$ is the foreign implementation that provides a way to evaluate $f$ to a
graph and value (as required in~\secref{core:foreign}). This general pattern of combining subgraphs also
applies to evaluated subterms in other contexts; for example the \ruleName{$\evalS$-let} rule for
$\exLet{x}{e}{e'}$ simply unions the graphs of subterms $e$ and~$e'$, and similarly for
\ruleName{$\evalS$-let-rec} and \ruleName{$\evalS$-app}.

To explain the last two rules, \ruleName{$\evalS$-let-rec} and  \ruleName{$\evalS$-app}, we first explain how
we support mutual recursion. \figref{core:close-defs} defines the auxiliary relation $\gamma, \rho, V
\closeDefs \gamma', G$, which takes a set $\rho$ of recursive definitions and produces the corresponding
environment $\gamma'$ of closures, along with a dependence graph $G$. This works by turning each function
definition $\rho(x_i) \in \rho$ into a newly allocated closure $(\gamma, \rho, \rho(x_i))_{\alpha_i}$ with
fresh address ${\alpha_i}$, capturing $\rho$ and a copy of the original environment $\gamma$, where each fresh
closure depends on ambient inputs $V$. The final dependence graph is the union of the graphs associated with
each closure.

Returning to \figref{core:eval}, rule \ruleName{$\evalS$-let-rec}  for  $\exLetRec{\seq{\bind{x}{\sigma}}}{e}$
is then the same as regular let-bindings, except using $\closeDefs$ to build the set of recursive definitions
$\seq{\bind{x}{\sigma}}$. The final rule, \ruleName{$\evalS$-app} for $\exApp{e}{e'}$, is especially important
to the graph semantics, as it is during application that we grow the ambient inputs $V$ that a given
evaluation step depends on \todo{Not quite, application is actually where we discard the current $V$ and
transition to a new active function context}. Here, we compute the closure
$\annClosure{\gamma_1}{\rho}{\sigma}{\alpha}$ of function $e$, and use $\sigma$ to pattern-match against the
value $v'$ of argument $e'$. Pattern-matching returns the selected branch $e^\twoPrime$ and the vertex set
$V'$ representsing the partial values in $v'$ consumed by pattern-matching in order to select $e^\twoPrime$.
The last premise, which evalutes the branch, then specifies its ambient dependencies as $V' \disjunion
\set{\alpha}$; that is, the (closure of the) function that was applied, $\alpha$, and the (consumed values of
the) argument it was applied to, $V'$.

% \subsubsection{Products of Projections of Views}

% Fix $A,B, C,D$ to range over Boolean algebras.
% \newline
% \begin{salign}
%    \textrm{Meet and Join} \\
%    \joinF &: X \to X \times X & \meetF &: X \times X \to X \\
%    \bwd{\joinF}(a,a') &= a \vee a' & \bwd{\meetF}(a) &= (a, a) \\
%    \fwd{\joinF}(a)    &= (a, a) & \fwd{\meetF}(a, a') &= a \wedge a'
% \end{salign}
% \begin{salign}
%    \textrm{Projections/Injections} \\
%    \inFst &: A \to A \times B & \inSnd &: B \to A \times B \\
%    \bwd{\inFst}(a, b) &= a & \bwd{\inSnd}(a, b) &= b \\
%    \fwd{\inFst}(a) &= (a, \top) & \fwd{\inSnd}(b) &= (\top, b)\\
%    \dual{\inFst} &: A \times B \to A & \dual{\inSnd} &: A \times B \to B \\
%    \bwd{\dual{\inFst}}(a) &= (a, \bot) & \bwd{\dual{\inSnd}}(b) &= (\bot, b) \\
%    \fwd{\dual{\inFst}}(a, b) &= a & \fwd{\dual{\inSnd}}(a, b) &= b
% \end{salign}
% Assume $f: A \to C$, $g: A \to D$
% \begin{salign}
%    \textrm{Products} \\
%    \funProd{f}{g} &: A \to C \times D \\
%    \bwd{\funProd{f}{g}}(c, d) &=\bwd{ f}(c) \vee\bwd{ g}(d) \\
%    \fwd{\funProd{f}{g}}(a) &= (\fwd{f}(a)\fwd{,g}(a)) \\
%    \dual{\funProd{f}{g}}&: C \times D \to A \\
%    \bwd{\dual{\funProd{f}{g}}}(a) &= (\bwd{f}(a),\bwd{ g}(a)) \\
%    \fwd{\dual{\funProd{f}{g}}}(c, d) &=\fwd{ }(c) \wedge\fwd{ g}(d)
% \end{salign}

% \begin{figure}
%    \begin{tikzcd}
%       C \arrow[rr, "\bwd{\meetF}"]   &  & C \times C \arrow[dd, "\bwd{f}"', shift right=2] \arrow[dd, "\opName{id}", shift left=2] \\
%                                   &  &                                                                                          \\
%    A \arrow[uu, "\dual{\fwd{f}}"] &  & A \times C
%    \end{tikzcd}
%    \caption{linked inputs with a single output view}
% \end{figure}
